% 01 — Cover and abstract
\thispagestyle{cover}
\vspace*{5mm}
\noindent{\sffamily\small\color{teal}\bfseries EFFECTIVE AUTO RESEARCH}\par
\vspace{2mm}
\noindent{\sffamily\fontsize{58}{62}\selectfont\bfseries EAR}\par
\vspace{3mm}
\noindent{\sffamily\fontsize{25}{33}\selectfont\bfseries 以人时为核心，\par\noindent 加速科研全生命周期}\par
\vspace{4mm}
\noindent{\sffamily\large\color{muted}问题发现 · 自主数学研究 · 评审回应}\par
\vspace{4mm}
\ReportAuthors
\vspace{4mm}
\noindent\textcolor{teal}{\rule{\textwidth}{1.2pt}}\par
\vspace{3mm}
\Lead{把人的时间留给关键研究判断。}
科研需要原创判断，也需要反复检索、组织输入、推进验证和修订稿件。EAR 将这些工作组织成可持续执行、可依据反馈改进的研究流程。第一目标是在质量要求下减少获得有效成果所需的人时；研究周期与机器费用提供辅助尺度。

问题发现模块以检索、批判、评分和剪枝形成可执行提案。数学引擎的内循环推进完整研究，外层遗传算法从研究结果中改进可复用策略。评审回应模块从已有证据和策略出发，用代理反馈集中解决决定回应质量的缺口。人的角色贯穿目标设定、关键技术验收和最终责任。

本报告给出机制、历史聚合结果和可重算的发布材料。统计截至 G3：固定的 28 个完成稿件中，当前轮模型评阅通过数从初轮 3 个增至第三轮修订后的 12 个。G3 所选策略的通过率为 75\%（$3/4$），基线为 50\%（$2/4$）；每个通过结果的研究调用减少 52.8\%。这些结果分别支持研究修订与策略演化的观察结论。\src{1,10}
\vspace{3mm}
\begin{insight}
\textbf{一条共同主线}\quad 尽早判断什么值得做，持续推进有反馈的研究，将已有成果转化为有证据的回应。Agent 接手反复执行，人保留少量高价值判断。
\end{insight}
\vfill
\noindent{\small\color{muted}中文技术报告 · 2026 年 10 月 · 公开聚合数据与可重建源码}

% 02 — First principles
\Page{01 / 第一性目标}{让有限人时用于高价值判断}{以研究质量为前提，减少主动投入；周期与机器费用分别记录。}
\Lead{科研自动化的价值，取决于完成可信研究需要多少人的主动工作。}
EAR 用以下目标组织系统设计：
\begin{equation}
\min H\qquad\text{subject to}\qquad Q\ge Q_{\min},\quad C\le B.
\label{eq:human-objective}
\end{equation}
$H$ 是主动人时，$Q_{\min}$ 是约定的研究质量，$C$ 是模型、工具与算力费用，$B$ 是预算。日历周期 $T$ 单独报告，必要时作为截止时间条件。机器并行可以缩短周期；人的主动投入要从实际工作记录计算。

\Sub{人时要覆盖完整工作，而非只看生成速度}
阅读检索、问题比较、推导实现、任务编排、结果检查、模型错误修复和失败路线返工都计入 $H$。多人协作时，人时为各人的主动投入之和。无人值守的运行等待计入周期，在等待期间没有主动工作的部分不计入人时。因此，委派执行带来的节省与新增审核负担需要在同一任务里结算。

\Sub{从目标推出三个环节}
\textbf{上游尽早筛选。}选错题会使后续推导、实现和写作一起失去价值。检索和批判应在大规模投入前提供继续、改写或退出的依据。

\textbf{中游持续执行。}保存设定、局部结果、失败路线与下一步，能够让 Agent 在已有状态上推进研究，减少人反复恢复背景和逐条下指令的工作。

\textbf{下游围绕证据修订。}评阅暴露的关键疑虑应回到证明、代码或实验。把主张与证据对应起来，能使人工检查集中到决定性技术位置。

\begin{insight}
\textbf{怎样判断有没有帮助}\quad 后文用检索记录、稿件评分和调用次数，说明系统目前做到了什么。它们对应减少重复工作的设计目标；实际节省多少人时，还需要在使用中记录。
\end{insight}

% 03 — Overview and evidence map
\Page{02 / 系统总览}{三个环节接起来，人保留关键判断}{目标、工具、持久工件与反馈，构成共同的研究接口。}
\Fig{\resizebox{.96\textwidth}{!}{\SystemOverview}}{EAR 的研究生命周期。提案、研究工件与评阅材料连接三个环节，各模块也支持独立输入和使用。}
\begin{center}\small
\begin{tabularx}{\textwidth}{P{27mm}Y Y}
\toprule
\textbf{模块} & \textbf{输入与产物} & \textbf{人保留的判断}\\
\midrule
问题发现 & 方向与条件 $\to$ 文献比较、候选、提案 & 问题价值与执行路径\\
数学研究 & 方向或提案 $\to$ 论证、代码、稿件与评价 & 关键技术结论与验收\\
评审回应 & 论文、评审与证据 $\to$ 回应及证据账本 & 主张边界与最终责任\\
\bottomrule
\end{tabularx}
\end{center}
\captionof{table}{三类模块的工件接口与人工判断。输入明确后，Agent 承担可委派的持续工作。}

\Sub{四项主张，对应四类证据}
\textbf{C1：证据驱动的候选修订与剪枝。}E1 展示新增近似工作怎样改变候选分数和继续建议。

\textbf{C2：内循环改善固定稿件子集的观察评阅指标。}E2 用同一完成稿件集合跟踪各轮平均评分与通过数。

\textbf{C3：演化策略改善历史选择比较中的通过率与调用产出。}E3 展示 G0 初始化与 G1--G3 三代演化，给出该阶段的种群走势、选择批次与成本分解。

\textbf{C4：实现有证据的回应优化，并提供反馈评价器证据。}E4 连接可执行策略循环与公开评阅数据上的小样本评价。\src{1}

% 04 — Idea search mechanics
\Page{03 / 问题发现}{批判性树搜索，让问题逐轮具体化}{每次扩展后重新检索、批判和评分，再决定下一次投入。}
\Lead{候选节点逐步明确问题、假设、方法与验证计划；新证据可以改变整条研究路线。}
\Fig{\resizebox{.94\textwidth}{!}{\CriticalTree}}{检索、候选扩展、多视角批判与剪枝构成循环。评分进入每轮选择，碰撞、不可行或重复的分支保留退出理由。}
\begin{center}\small
\begin{tabularx}{\textwidth}{P{28mm}Y P{16mm}}
\toprule
\textbf{评价字段} & \textbf{它为什么影响投入} & \textbf{权重}\\
\midrule
Novelty & 检查最近工作与核心机制的实质差异 & .25\\
Venue & 评议研究价值、严谨性与目标要求 & .35\\
Strategic & 检查可证伪性、证据可得性与资源适配 & .20\\
Feasibility & 检查实现复杂度、算力、数据与时间线 & .20\\
\bottomrule
\end{tabularx}
\end{center}
\captionof{table}{当前四类模型评价字段与组合权重。评价理由与检索来源一并保存。}

当前实现用 UCT 引导 best-first 扩展：模型生成子节点，工具更新访问次数与价值，并沿父链回传。每个新候选至少重新进行新颖性和研究适配检查。评分与可核查的碰撞、验证规则和结构信号共同指导投入；研究点子不使用随机 rollout 模拟终局。多 Agent 表示任务职责分工，多评阅者画像的实际调用按对应提示与记录解释。\src{1,2}

% 05 — Harness integration
\Page{04 / 技术接入}{把研究意图变成调用、反馈与持久工件}{人通过命令交互，Claude Code Harness 调度，Codex MCP 批判评分。}
\Lead{研究策略只有落到实际输入、工具调用和产物上，才能在下一步继续使用。}
\Fig{\resizebox{.96\textwidth}{!}{\HarnessAccess}}{命令文件组织调研、生成、筛选和精修；Codex MCP 返回批判与评分，Claude Code 处理反馈并保存报告。}
\Sub{一次交互怎样推进研究？}
人用 \texttt{/start} 给方向与资源条件，也可单独执行 \texttt{/idea-screen} 或 \texttt{/idea-refine}。命令文件指定任务、输入、工具和期望产物。筛选将候选、核心主张和近邻文献交给 Codex MCP，要求给出新颖性、差异、继续或放弃建议，以及判断理由。Harness 接收反馈并保存候选排序与精修材料，下一次交互直接使用这些工件。\src{1}

\Sub{为什么需要持久报告？}
会话中的一句建议难以支持协作与复核。报告保存最近工作、候选差异、问题锚点和验证计划，人可以检查关键取舍，新的执行者也能够从共同设定开始。调用记录还区分了实际工具评价与宿主自审，使反馈来源可追踪。

\begin{insight}
\textbf{执行接口的作用}\quad 人提供目标与关键追问；命令和 Harness 接手组织输入、调用评价、落盘和后续衔接。早期接入采用交互式命令流程，后续树搜索进一步组织候选扩展，两种实现沿同一研究接口演进。
\end{insight}

% 06 — E1 idea evidence
\Page{05 / E1：证据改变选择}{新增近似工作，使高分候选退出}{筛选应当改变投入决定，而不只是给想法增加一个分数。}
\Lead{一次复判在两个候选周围找到三项近似工作，其中一个候选的新颖性由 $8/10$ 降至 $4/10$，建议转为放弃。}
\begin{insight}
\textbf{证据 $\to$ 判断 $\to$ 行动}\\[1mm]
首轮候选与初始筛选\quad $\longrightarrow$\quad 二轮补充近似工作\\[1mm]
重新比较核心差异\quad $\longrightarrow$\quad $8\to4$\quad $\longrightarrow$\quad 放弃该分支
\end{insight}

\Sub{为什么第二轮会改变结论？}
初始检索没有充分覆盖最近工作，原先支持候选的差异因此显得更大。继续检索找到近似机制后，系统将新增来源与核心主张逐项比较，发现原建议缺少足够的差异支撑。分数更新不是单独发生的数字变化，而是对应文献、差异分析和继续建议的同时修订。

\Sub{退出理由也是研究产物}
树中保留被剪枝的节点、碰撞依据与判断理由。研究者可以复核为什么停止；后续搜索可以利用这些信息，避免重新投入已经缺乏差异的路线。其他候选仍可依据自己的证据继续推进，三项近似工作并非全部对应这个被放弃的候选。

\Sub{从候选列表到可执行提案}
对留下的方案，精修冻结问题锚点，明确假设、主要论据、实现接口与最低验证计划。这样，人的阅读从广泛检索和反复比较转向少量关键候选的价值与可执行性。E1 将新增证据、重新判断与退出行动连接起来，使筛选可以直接影响后续执行。\src{1,10}

\begin{insight}
\textbf{一个完整搜索结果}\quad 既提供值得继续的可执行提案，也保留退出路线的来源与理由。检索产生的负向证据同样帮助研究者决定人的时间应投入在哪里。
\end{insight}

% 07 — Math flow and dual loop
\Page{06 / 数学研究}{内循环做研究，外循环改进研究做法}{完整研究进入固定评价，反馈回到可复用策略。}
\Lead{数学与算法任务有明确对象和可定位的反馈，适合把推理、实现、检查和修正连接成连续小循环。}
\Fig{\resizebox{.95\textwidth}{!}{\MathDualLoop}}{完整研究置于策略演化内部。内层推进当前任务，外层比较策略、保留精英并生成下一代；固定模型评阅提供共同参照。}
\Sub{内层为什么能够持续？}
精确命题、失败的证明步骤、反例、数值检查和实验对照为下一步提供依据。模型按状态选择行动，保存设定、局部结果与失败路线，修正后继续推进。稿件评阅指出的问题可以返回技术工作，也可以触发表述修订，直至达到交付条件或记录退出。

\Sub{外层为什么需要完整研究结果？}
提示看起来合理，不等于能产生更好的研究。外层让每个策略运行完整内层，观察目标选择、技术尝试、稿件与评阅结果，再根据失败画像改进做法。已有 Harness 提供初始种子，精英、变异、交叉与探索构成后续种群。

\begin{insight}
\textbf{两个循环共享一个目标}\quad 内层接手当前研究的反复执行，外层将反馈转为后续可复用策略。人保留研究方向、关键技术验收与最终责任，减少反复切工具、重建上下文和逐步编排的需求。\src{1,3,6,7}
\end{insight}

% 08 — Evolution technical detail
\Page{07 / 遗传调优}{改的是研究策略，评价来自完整研究}{策略基因、演化算子与固定评阅，使改进能够被后续执行检验。}
\Sub{可修改范围与共同参照}
历史实验的基因组是阶段提示与参数，改变目标筛选、证明组织、写作和修订方式。当前便携实现将接口扩展到工作代码、提示、技能和配置。候选研究策略与初始策略 分开保存；固定任务列表、模型评阅及标准置于候选可修改范围之外。每次改动都运行完整研究，并留下谱系、结果和失败画像。\src{1}

\Sub{反馈怎样变成下一代？}
\textbf{精英保留}保存表现较好的策略；\textbf{定向变异}根据退出原因、技术缺口或低分评阅修改父策略；\textbf{语义交叉}协调两个父策略的互补指令；\textbf{探索}允许较大变化。演化发生在研究方法上，不改变基础模型权重。历史改动包括收窄候选清单、提高目标成功阈值、在选题前要求明确核心引理与可行证明路线。

\Sub{历史适应度与当前接口}
历史 episode 的适应度为：
\[
F(e)=\begin{cases}
q_{\mathrm{best}}(e),&\text{完成稿件且具有有效评分},\\
0.5,&\text{研究退出},\\
0,&\text{筛选淘汰},\\
-0.5,&\text{没有目标或搜索失败},\\
-2,&\text{预算耗尽或贡献过于平凡}.
\end{cases}
\]
其中 $q_{\mathrm{best}}$ 是评分与修订后保留的最佳稿件均分。该轮选择按平均 $F$ 排序；历史精英档案则等权平均归一化 $F$、完成率与模型评阅通过率。两项选择各有记录；完成稿件但所有评分失败时，历史实现使用 3.0 回退值。当前接口等权结合归一化质量、完成率和固定终审通过率，相当时再比较调用与返工。效率结果按真实调用计数另行报告。反思式提示演化同样强调从失败反馈中改进可复用策略。\src{1,4,10}

% 09 — Evaluation setup
\Page{08 / 数学评测设定}{先固定单位和口径，再比较修订与演化}{展示 G0 初始化与 G1--G3 三代演化，保留范围内的全部尝试。}
\Lead{G0--G3 共 64 个研究 episode，记录 630 次研究调用、28 个完成稿件和 15 个 selected-best 模型评阅通过结果。}
每个策略批次执行四个完整 episode。策略从给定方向中选择具体目标，因而比较同时覆盖选题与研究做法。初始策略 在每代保留，候选策略按该轮观察表现选择。固定的 28 个完成稿件构成内循环逐轮比较的子集；外层比较使用每代策略批次及其新研究结果。

\begin{center}\small
\begin{tabularx}{\textwidth}{P{34mm}Y}
\toprule
\textbf{单位或定义} & \textbf{本报告中的含义}\\
\midrule
研究 episode & 一次完整研究尝试，失败也计入总数与调用\\
研究调用 & 历史研究驱动器记录的阶段调用次数\\
完成稿件 & 完成所需稿件产物的研究结果\\
当前轮通过 & 该评分轮的模型评阅聚合标签为 Accept\\
Selected-best 通过 & 最高均分版本恢复后，对应的评阅标签为 Accept\\
固定内循环子集 & 同一组 28 个完成稿件，跟踪初轮与三轮修订\\
\bottomrule
\end{tabularx}
\end{center}
\captionof{table}{数学评测单位。完成、当前轮通过与最佳版本通过具有不同含义，分别用于对应分析。}

\Sub{历史评阅与版本恢复}
每轮由三个新模型会话评价稿件，至少一票 Accept 即记为通过。初评之后执行三轮修订与再次评阅，结束时交付均分最高的稿件版本，并采用该版本的评阅结果。若把要求改为三票均 Accept 且每票至少 7 分，G3 两组各有 $1/4$ 次尝试满足要求。下文的通过数均采用前述历史规则。

本报告从已有记录中选取 G0--G3 展示三代演化过程，范围在查看记录后确定；所有统计都对应这一阶段。完整评分规则与成本计算保存在随报告发布的数据中。\src{10}

% 10 — Inner results
\Page{09 / E2：内循环}{固定稿件子集的评阅指标逐轮改善}{同一组 28 个完成稿件，分别观察当前轮均分与通过数。}
\Fig{\earfigure{inner_progress}}{G0--G3 的 28 个完成稿件。当前轮通过数依次为 3、10、10、12；均分依次为 4.647500、5.247143、5.497857、5.674643。}

\Sub{修订回到具体缺口}
内循环把模型评阅发现的问题转为研究行动：技术缺口触发推导或实现修复，贡献不足促使范围与主张重新审查，表述问题进入稿件重写。再次评阅提供新反馈，保留完整轨迹，使人能定位哪些内容被修改以及质量判断怎样变化。

在固定子集中，当前轮通过比例由 $3/28=10.7\%$ 增至 $12/28=42.9\%$，平均评分由约 4.65 增至 5.67。固定分母避免将更多稿件进入集合误读为单稿改进。这一结果支持 C2：研究修订与该完成稿件子集的观察评阅指标改善相联系。

\begin{insight}
\textbf{两种结果，各自回答一个问题}\quad 当前轮轨迹展示修订过程；最终 selected-best 的 15 个通过结果展示交付时保留的版本。它们来自不同统计口径，最佳版本恢复也保留了最后一轮回落时的较好稿件。
\end{insight}

% 11 — Evolution results
\Page{10 / E3：外层演化}{三代演化后，选中策略的通过率达到 75\%}{每一代都用完整研究尝试，检验修改后的策略有没有帮助。}
\Fig{\earfigure{evolution_progress}}{每代选中策略的模型评阅通过率。G0 是初始化，G1--G3 是三代演化；每个点统计四次研究尝试，包含失败和未完成的尝试。}

\Sub{把上一轮遇到的问题，变成下一轮的改法}
每一代结束后，系统整理失败原因和评阅意见，修改选题要求、证明组织和修订提示，再运行一批研究。记录中的改动包括缩短候选清单、在选题前要求明确核心引理和证明路线，以及提高贡献要求。这些改动通过后续研究结果来检验。

\Sub{这三代里，通过率从 0\% 提高到 75\%}
G0 和 G1 选中策略的四次尝试均未通过；G2 有两次通过，G3 有三次通过。图中的 75\% 因此对应 $3/4$ 次尝试。每个点都使用四次尝试作为分母，能够直接看到选中策略在这一阶段的变化。

图中展示的是每代选中的策略。计入 G3 所有候选策略与初始策略后，整代共有 20 次尝试，9 次通过，通过率为 45\%。各代完整计数保存在随报告发布的数据中。\src{10}

\begin{insight}
\textbf{外循环改进的是研究做法}\quad 把一次研究暴露的问题写进下一版策略，让后续尝试用上已经得到的经验。研究者可以查看修改理由和结果，再决定是否采用这套做法。
\end{insight}

% 12 — Cost decomposition
\Page{11 / E3：调用与产出}{每个通过结果，研究调用减少 52.8\%}{比较 G3 的初始策略与选中策略，分别计算研究、评阅和修订调用。}
\Fig{\resizebox{.84\textwidth}{!}{\earfigure{cost_comparison}}}{G3：研究调用每通过由 $41/2=20.50$ 降至 $29/3=9.67$，减少 52.8\%；研究、评阅与修订合计由 $86/2=43.00$ 降至 $74/3=24.67$，减少 42.6\%。}
\begin{center}\small\setlength{\tabcolsep}{3pt}
\begin{tabularx}{\textwidth}{P{11mm}P{21mm}rrrrrr}
\toprule
\textbf{批次} & \textbf{策略} & \textbf{研究} & \textbf{评阅} & \textbf{修订$^*$} & \textbf{合计} & \textbf{通过} & \textbf{合计/通过}\\
\midrule
G3 & 初始策略 & 41 & 36 & 9 & 86 & 2 & 43.00\\
G3 & 选择 & 29 & 36 & 9 & 74 & 3 & 24.67\\
\bottomrule
\end{tabularx}
\end{center}
\captionof{table}{G3 选择批次的调用账本，每组四次研究。评阅是记录的评阅者调用；修订$^*$按非初轮推得。}

G3 研究调用每 episode 从 10.25 降至 7.25，同时通过比例由 $2/4$ 增至 $3/4$，共同带来每通过调用下降。G0--G3 全部策略累计 630 次研究调用，产生 15 个通过结果，即累计通过率 $15/64=23.4\%$、每通过 42 次研究调用。这一总账包含策略搜索投入，与只使用所选策略的成本分别报告。代理不含完整教练调用、内部工具与计费 token，覆盖范围随协议一并发布。

% Rebuttal and closing pages, revised from a first-time reader's perspective.
\Page{12 / 评审回应}{收到评审后，先处理最关键的问题}{找到已有答案，只为缺口补研究，再组织回应。}
\Lead{作者首先要判断：哪些疑虑能用现有证明和实验回答，哪些确实需要补做研究。}
\Fig{\resizebox{.95\textwidth}{!}{\RebuttalFlow}}{从论文和评审出发，合并相同问题，找到对应证据。需要补研究时先补齐，再起草、检查和修改回应。}

\Sub{先找答案，再决定补什么}
同一个问题可能被几位审稿人反复提出。EAR 先把这些意见放在一起，找到论文中的对应论据，再判断缺的是解释、证明还是实验。已有材料能够回答的，就组织清楚；确实缺证据的，才安排补充研究。同一项补充实验可以用于回答多条相关意见。\src{1,5,9}

\Sub{把时间花在最影响评价的问题上}
社区公开的回应方法和作者经验提供起点。系统比较不同的回答顺序、证据选择和写法，优先处理影响论文判断的疑虑。实现时，将这个任务写成一个优化目标：
\[
\max_{\theta}\ \mathbb E[\Delta R\mid\mathrm{Context},\theta]
\quad\text{subject to}\quad H\le H_B,\ C\le B,\ T\le D.
\]
这里的 Context 就是论文、评审和已有证据；$\theta$ 是回应策略，$\Delta R$ 是希望改善的评分，$D$ 是截止时间。含义很直接：在可用的人时、费用和期限内，选择最有希望解决关键疑虑的回应方式。

\begin{insight}
\textbf{模型给建议，作者检查关键内容}\quad 模型检查问题有没有回答、论据有没有依据，并预测评分可能上涨、不变或下降。这些反馈用于修改草稿；新增技术结论和最终回应由作者检查。
\end{insight}

\Page{13 / 回应策略优化}{保留好草稿，继续解决剩下的问题}{从社区已有方法开始，每轮都接着上一轮的结果改。}
\Fig{\resizebox{.95\textwidth}{!}{\RebuttalPopulation}}{尝试不同写法，比较模型评价，保留较好的草稿。下一轮带上尚未回答的问题和修改建议，继续完善。}

\Sub{先尝试几种有依据的写法}
初始策略来自公开的回应方法和作者经验：可以先解释已有证据，也可以先回应审稿人最关心的问题，或先安排缺失的验证。它们使用同一份论文和评审材料，区别在于如何组织工作和回答。

\Sub{每轮留下草稿、问题和建议}
系统保留当前较好的回应，并记下哪里还没讲清、哪些证据还不足、下一步应该怎么改。下一轮直接接着这些材料推进。实现上，这是黑盒优化：通过比较候选回应和读取反馈来调整策略，无需计算评分对写作策略的导数。

\Sub{记录中确实发生了策略调整}
一份匿名两轮记录中，第一轮选中的草稿未达到内部模型的评价要求。系统保留草稿和修改建议，第二轮换用另一策略后通过评价。这展示了系统会根据反馈调整做法，记录中的结果来自内部模型检查。

草稿通过内容评价后，系统停止尝试本轮其余策略，再检查回应是否忠实于论文和证据、表述是否准确。发现问题就继续修，检查完成后交给作者审阅。\src{1,10}

\Page{14 / E4：回应效果预测}{模型能预测回应后的评分变化吗？}{用已经发生的评审案例，检查模型的判断有没有依据。}
\Lead{模型读取论文、初始评审和作者当时的回应，预测评分会上涨、不变还是下降。}
最终评分只用来核对预测。初始评分和最终评分的差，给出真实变化标签；每个案例单独输入模型，避免不同案例之间混入信息。测试结果如下。\src{1,8,10}

\begin{center}\small
\begin{tabularx}{\textwidth}{P{32mm}Y P{21mm}}
\toprule
\textbf{历史案例} & \textbf{预测任务} & \textbf{Macro-F1}\\
\midrule
36 例均衡子集 & 区分提分与不变，各 18 例 & .805\\
57 例完整测试集 & 区分提分、不变和降分 & .503\\
同一 57 例来源集 & 来源报告的提分／不变二类指标 & .755\\
\bottomrule
\end{tabularx}
\end{center}
\captionof{table}{DeepSeek V4 Pro 的记录结果。Macro-F1 分别计算每个类别的 F1 后取平均，取值越接近 1 越好。三行对应不同测试设置。}

\Sub{能区分提分与不变，识别降分仍然较弱}
在 36 例均衡测试中，模型区分提分和不变的 macro-F1 为 .805。57 例完整测试集中包含 6 例降分情况，三类 macro-F1 为 .503，其中降分类的 F1 为 0。这提示系统不能只看一个总体分数，还要检查模型在哪一类判断上容易出错。

\Sub{把反馈用于改草稿，也检查它的依据}
EAR 将评分变化预测与证据检查一起用于比较回应。测试中，进一步修改提示没有得到更好的留出集成绩，因此保留原始零样本提示，没有修改模型权重。

\begin{insight}
\textbf{这些数字回答的是预测能力}\quad 测试使用已经写好的历史回应。系统生成新回应后，真实审稿人是否因此提高评分，需要用另一项实验检验。
\end{insight}

\Page{15 / 对研究者意味着什么}{让研究持续推进，把关键判断留给人}{从筛选问题到修订论文，再到组织回应，EAR 接手反复执行的工作。}
研究会反复经历同一类过程：查一轮文献，发现缺口；做一次尝试，遇到问题；补完证据，再修改结论。每次重新整理背景、安排下一步和检查遗漏，都会占用研究者的时间。EAR 把这些步骤接起来，让已经得到的材料和反馈能够继续用于下一步。

\Sub{找问题：让新证据改变投入方向}
搜索模块持续比较候选与已有工作。记录中，新找到的近似工作使一个候选被放弃，并留下退出理由。研究者可以先检查这些关键差异，再决定是否投入更大规模的推导和实验。

\Sub{做研究：把尝试、检查和修订接起来}
内循环保留证明、代码、失败路线和评阅意见，按具体问题继续修。固定的 28 个稿件中，当前版本通过模型评阅的数量从初轮 3 个增至第三轮修订后的 12 个。外循环再修改研究策略：在展示到 G3 的阶段中，当代选中策略四次尝试有三次通过；每个通过结果的研究调用，比同在 G3 运行的初始策略少 52.8\%。\src{10}

\Sub{写回应：把精力放在尚未解决的疑虑上}
回应模块合并重复问题，查找已有依据，必要时安排补充研究，再比较和修改草稿。每一轮保留有效内容和下一步建议，作者可以直接检查关键证据、技术结论和最终表述。

\begin{insight}
\textbf{一次判断，带动下一段工作}\quad 人给方向、设标准、检查关键结果；Agent 接着检索、尝试和修订。研究者回来时，有新的结果和清楚的修改记录，可以据此决定下一步。
\end{insight}

\noindent 本报告展示了流程、模型评阅和调用结果。实际节省多少主动人时，还需要在使用中直接记录。

